\frfilename{mt231.tex} 
\versiondate{25.8.15} 
\copyrightdate{1994} 
      
\def\chaptername{Teorema Radon-Nikod\'ym} 
\def\sectionname{Fungsional aditif terhitung} 
      
\newsection{231} 
      
Saya mulai dengan suatu uraian abstrak tentang objek-objek yang, dalam 
keadaan yang sesuai, akan bersesuaian dengan integral tak tentu dari 
fungsi-fungsi terintegralkan umum. Dalam bagian ini saya menyajikan 
bagian-bagian teori yang tidak melibatkan ukuran, melainkan hanya suatu 
himpunan beserta aljabar-sigma subhimpunannya yang dipilih secara khusus. 
Konsep-konsep dasarnya ialah konsep fungsional `aditif hingga' dan `aditif 
terhitung', dan terdapat satu teorema penting, yakni `dekomposisi Hahn' (231E). 
      
\leader{231A}{Definisi} Misalkan $X$ suatu himpunan dan $\Sigma$ suatu 
aljabar subhimpunan dari $X$\cmmnt{ (136E)}.   Suatu fungsional 
$\nu:\Sigma\to\Bbb R$ disebut {\bf aditif hingga}, atau cukup disebut 
{\bf aditif}, jika $\nu(E\cup F)=\nu E+\nu F$ kapan pun $E$, 
$F\in\Sigma$ dan $E\cap F=\emptyset$. 
      
\leader{231B}{Fakta elementer} Misalkan $X$ suatu himpunan, $\Sigma$ suatu 
aljabar subhimpunan dari $X$, dan $\nu:\Sigma\to\Bbb R$ suatu fungsional 
aditif hingga. 
      
\header{231Ba}{\bf (a)} $\nu\emptyset=0$.   \prooflet{(Sebab 
$\nu\emptyset=\nu(\emptyset\cup\emptyset)=\nu\emptyset+\nu\emptyset$.) 
} 
      
\header{231Bb}{\bf (b)} Jika $E_0,\ldots,E_n$ adalah anggota-anggota 
$\Sigma$ yang saling lepas, maka 
$\nu(\bigcup_{i\le n}E_i)=\sum_{i=0}^n\nu E_i$. 
      
\header{231Bc}{\bf (c)} Jika $E$, $F\in\Sigma$ dan $E\subseteq F$, maka 
$\nu F=\nu E + \nu(F\setminus E)$.   Secara lebih umum, untuk sembarang 
$E$, $F\in\Sigma$, 
      
\Centerline{$\nu F=\nu(F\cap E)+\nu(F\setminus E)$,}

\Centerline{$\nu E+\nu F
\prooflet{\mskip5mu=\nu E+\nu(F\setminus E)+\nu(F\cap E)}
=\nu(E\cup F)+\nu(E\cap F)$,}

\Centerline{$\nu E-\nu F 
\prooflet{\mskip5mu = \nu(E\setminus F)+\nu(E\cap F) 
-\nu(E\cap F)-\nu(F\setminus E)} 
=\nu(E\setminus F)-\nu(F\setminus E)$.} 
      
\leader{231C}{Definisi} Misalkan $X$ suatu himpunan dan $\Sigma$ suatu 
aljabar subhimpunan dari $X$.   Suatu fungsi 
$\nu:\Sigma\to\Bbb R$ disebut {\bf aditif terhitung} atau 
{\bf $\sigma$-aditif} jika $\sum_{n=0}^{\infty}\nu E_n$ ada di $\Bbb 
R$ dan sama dengan $\nu(\bigcup_{n\in\Bbb N}E_n)$ untuk setiap 
barisan saling lepas $\sequencen{E_n}$ di dalam $\Sigma$ sedemikian sehingga 
$\bigcup_{n\in\Bbb N}E_n\in\Sigma$. 
      
\cmmnt{\medskip 
      
\noindent{\bf Catatan} Perhatikan bahwa ketika saya menggunakan frasa 
`fungsional aditif terhitung', saya bermaksud meniadakan kemungkinan 
$\infty$ sebagai nilai fungsional tersebut.   Dengan demikian, suatu ukuran 
merupakan fungsional aditif terhitung jika dan hanya jika ukuran tersebut 
berhingga total (211C). 
      
Kadang-kadang Anda akan melihat frasa `{\bf ukuran bertanda}' digunakan 
untuk menyatakan apa yang saya sebut fungsional aditif terhitung. 
} 
      
\leader{231D}{Fakta elementer} Misalkan $X$ suatu himpunan, $\Sigma$ suatu 
aljabar-sigma subhimpunan dari $X$, dan $\nu:\Sigma\to\Bbb R$ suatu 
fungsional aditif terhitung. 
      
\header{231Da}{\bf (a)} $\nu$ bersifat aditif hingga.   \prooflet{\Prf\ 
(i) Dengan menetapkan $E_n=\emptyset$ untuk setiap $n\in\Bbb N$, 
$\sum_{n=0}^{\infty}\nu\emptyset$ harus terdefinisi di $\Bbb R$, sehingga 
$\nu\emptyset$ harus sama dengan $0$.   (ii) Sekarang, jika $E$, 
$F\in\Sigma$ dan $E\cap F=\emptyset$, kita dapat menetapkan $E_0=E$, 
$E_1=F$, $E_n=\emptyset$ untuk $n\ge 2$ dan memperoleh 
      
\Centerline{$\nu(E\cup F)=\nu(\bigcup_{n\in\Bbb N}E_n) 
=\sum_{n=0}^{\infty}\nu E_n=\nu E+\nu F$.  \Qed} 
} 
      
\header{231Db}{\bf (b)} Jika $\sequencen{E_n}$ adalah barisan tak-menurun 
di dalam $\Sigma$, dengan gabungan $E\in\Sigma$, maka 
      
\Centerline{$\nu E 
=\nu E_0+\sum_{n=0}^{\infty}\nu(E_{n+1}\setminus E_n) 
=\lim_{n\to\infty}\nu E_n$.} 
      
\header{231Dc}{\bf (c)} Jika $\sequencen{E_n}$ adalah barisan tak-meningkat 
di dalam $\Sigma$ dengan irisan $E\in\Sigma$, maka 
      
\Centerline{$\nu E 
\prooflet{\mskip5mu =\nu E_0-\lim_{n\to\infty}\nu(E_0\setminus E_n)} 
=\lim_{n\to\infty}\nu E_n$.} 
      
\header{231Dd}{\bf (d)} Jika $\nuprime:\Sigma\to\Bbb R$ adalah fungsional 
aditif terhitung lainnya, dan $c\in\Bbb R$, maka 
$\nu+\nuprime:\Sigma\to\Bbb R$ dan $c\nu:\Sigma\to\Bbb R$ bersifat 
aditif terhitung. 
      
\header{231De}{\bf (e)} Jika $H\in\Sigma$, maka 
$\nu_H:\Sigma\to\Bbb R$ bersifat aditif terhitung, dengan 
$\nu_HE=\nu(E\cap H)$ untuk setiap 
$E\in\Sigma$.   \prooflet{\Prf\ Jika $\sequencen{E_n}$ adalah barisan 
saling lepas di dalam 
$\Sigma$ dengan gabungan $E\in\Sigma$, maka $\sequencen{E_n\cap H}$ juga 
saling lepas, dengan gabungan $E\cap H$, sehingga 
      
\Centerline{$\nu_HE 
=\nu(H\cap E) 
=\nu(\bigcup_{n\in\Bbb N}(H\cap E_n)) 
=\sum_{n=0}^{\infty}\nu(H\cap E_n) 
=\sum_{n=0}^{\infty}\nu_HE_n$.  \Qed}} 
      
\cmmnt{\medskip 
      
\noindent{\bf Catatan} Untuk sementara, kita akan menggunakan konsep 
`fungsional aditif terhitung' hanya pada aljabar-sigma $\Sigma$, 
sehingga kita dapat menganggap bahwa gabungan dan irisan di atas termasuk 
dalam $\Sigma$. 
}%end of comment 
      
\leader{231E}{}\cmmnt{ Semua gagasan di atas hanyalah modifikasi kecil 
dari gagasan-gagasan yang sudah diperlukan pada awal sekali teori ruang 
ukur.   Sekarang kita sampai pada sesuatu yang lebih berbobot. 
      
\medskip 
      
\noindent}{\bf Teorema} Misalkan $X$ suatu himpunan, $\Sigma$ suatu 
aljabar-sigma subhimpunan dari $X$, dan $\nu:\Sigma\to\Bbb R$ suatu 
fungsional aditif terhitung.   Maka 
      
(a) $\nu$ terbatas; 
      
(b) terdapat suatu himpunan $H\in\Sigma$ sedemikian sehingga 
      
\Centerline{$\nu F\ge 0$ kapan pun 
$F\in\Sigma$ dan $F\subseteq H$,} 
      
\Centerline{$\nu F\le 0$ kapan pun $F\in\Sigma$ dan 
$F\cap H=\emptyset$.} 
      
      
\proof{{\bf (a)} \Quer\ Andaikan, jika mungkin, sebaliknya.    Untuk 
$E\in\Sigma$, tetapkan $M(E)=\sup\{|\nu F|:F\in\Sigma,\,F\subseteq E\}$; 
maka $M(X)=\infty$.   Selain itu, kapan pun $E_1$, $E_2$, $F\in\Sigma$ dan 
$F\subseteq E_1\cup E_2$, kita mempunyai 
      
\Centerline{$|\nu F|=|\nu(F\cap E_1)+\nu(F\setminus E_1)| 
\le|\nu(F\cap E_1)|+|\nu(F\setminus E_1)|\le M(E_1)+M(E_2)$,} 
      
\noindent sehingga $M(E_1\cup E_2)\le 
M(E_1)+M(E_2)$.   Pilih suatu barisan $\sequencen{E_n}$ di dalam $\Sigma$ 
sebagai berikut.   Tetapkan $E_0=X$.   Jika $M(E_n)=\infty$, dengan 
$n\in\Bbb N$, maka tentu terdapat suatu $F_n\subseteq E_n$ sedemikian sehingga 
$|\nu F_n|\ge 1+|\nu E_n|$, sehingga 
$|\nu(E_n\setminus F_n)|\ge 1$.   Sekarang 
sekurang-kurangnya salah satu dari $M(F_n)$ dan $M(E_n\setminus F_n)$ 
tak berhingga; jika $M(F_n)=\infty$, tetapkan $E_{n+1}=F_n$; jika tidak, 
tetapkan $E_{n+1}=E_n\setminus F_n$; dalam kedua kasus, perhatikan bahwa 
$|\nu(E_n\setminus E_{n+1})|\ge 1$ dan $M(E_{n+1})=\infty$, sehingga 
induksi dapat dilanjutkan. 
      
Setelah menyelesaikan induksi ini, tetapkan 
$G_n=E_n\setminus E_{n+1}$ untuk $n\in\Bbb N$.   Maka 
$\sequencen{G_n}$ adalah barisan saling lepas di dalam 
$\Sigma$, sehingga $\sum_{n=0}^{\infty}\nu G_n$ terdefinisi di $\Bbb R$ dan 
$\lim_{n\to\infty}\nu G_n=0$; tetapi $|\nu G_n|\ge 1$ untuk setiap $n$. 
\Bang 
      
\medskip 
      
{\bf (b)(i)} Menurut (a), 
$\gamma=\sup\{\nu E:E\in\Sigma\}<\infty$.   Pilih suatu barisan 
$\sequencen{E_n}$ di dalam $\Sigma$ sedemikian sehingga 
$\nu E_n\ge\gamma-2^{-n}$ untuk setiap $n\in\Bbb N$.   Untuk 
$m\le n\in\Bbb N$, tetapkan $F_{mn}=\bigcap_{m\le i\le n}E_i$.   Maka 
$\nu F_{mn}\ge\gamma-2\cdot 2^{-m}+2^{-n}$ untuk setiap $n\ge m$.   \Prf\ 
Lakukan induksi pada $n$. Untuk $n=m$, pernyataan ini langsung diperoleh 
dari pemilihan $E_m=F_{mm}$. 
Untuk langkah induksi, kita mempunyai $F_{m,n+1}=F_{mn}\cap E_{n+1}$, 
sedangkan tentu saja $\gamma\ge\nu(E_{n+1}\cup F_{mn})$, sehingga 
      
$$\eqalignno{\gamma+\nu F_{m,n+1} 
&\ge\nu(E_{n+1}\cup F_{mn})+\nu(E_{n+1}\cap F_{mn})\cr&
=\nu E_{n+1}+\nu F_{mn}
\Displaycause{231Bc}
\ge\gamma-2^{-n-1}+\gamma-2\cdot 2^{-m}+2^{-n}
\Displaycause{berdasarkan pemilihan $E_{n+1}$ dan hipotesis 
induksi} 
=2\gamma-2\cdot 2^{-m}+2^{-n-1}.
\cr}$$ 
      
\noindent Dengan mengurangkan $\gamma$ dari kedua ruas, 
$\nu F_{m,n+1}\ge\gamma-2\cdot 2^{-m}+2^{-n-1}$ dan induksi 
berlanjut. \Qed 
      
\medskip 
      
\quad{\bf (ii)} Untuk $m\in\Bbb N$, tetapkan 
      
\Centerline{$F_m=\bigcap_{n\ge m}F_{mn}=\bigcap_{n\ge m}E_n$.} 
      
\noindent Maka 
      
\Centerline{$\nu F_m 
=\lim_{n\to\infty}\nu F_{mn}\ge\gamma-2\cdot 2^{-m}$,} 
      
\noindent menurut 231Dc. 
Selanjutnya, $\sequence{m}{F_m}$ tak-menurun, sehingga dengan menetapkan 
$H=\bigcup_{m\in\Bbb N}F_m$ kita mempunyai 
      
\Centerline{$\nu H=\lim_{m\to\infty}\nu F_m\ge\gamma$;} 
      
\noindent karena $\nu H$ 
tentu lebih kecil daripada atau sama dengan $\gamma$, maka $\nu H=\gamma$. 
      
\medskip

\quad{\bf (iii)} Jika $F\in\Sigma$ dan $F\subseteq H$, maka 
      
\Centerline{$\nu H-\nu F=\nu(H\setminus 
F)\le\gamma=\nu H$,} 
      
\noindent sehingga $\nu F\ge 0$.   Jika $F\in\Sigma$ dan $F\cap 
H=\emptyset$, maka 
      
\Centerline{$\nu H+\nu F=\nu(H\cup F)\le\gamma=\nu H$} 
      
\noindent sehingga $\nu F\le 
0$.   Ini menyelesaikan pembuktian. 
}%end of proof of 231E 
      
      
\leader{231F}{Akibat} Misalkan $X$ suatu himpunan, $\Sigma$ suatu 
aljabar-sigma subhimpunan dari $X$, dan $\nu:\Sigma\to\Bbb R$ suatu 
fungsional aditif terhitung.   Maka $\nu$ dapat dinyatakan sebagai selisih 
dua ukuran berhingga total dengan domain $\Sigma$. 
      
\proof{ Ambil $H\in\Sigma$ sebagaimana diuraikan dalam 231Eb.   Tetapkan 
$\nu_1 E=\nu(E\cap H)$, $\nu_2 E=-\nu(E\setminus H)$ untuk $E\in\Sigma$. 
Maka, seperti dalam 231Dd-e, $\nu_1$ dan $\nu_2$ keduanya merupakan 
fungsional aditif terhitung pada $\Sigma$, dan tentu saja 
$\nu=\nu_1-\nu_2$.   Namun, berdasarkan pemilihan $H$, $\nu_1$ dan 
$\nu_2$ juga keduanya nonnegatif, sehingga merupakan ukuran berhingga total. 
}%end of proof of 231F 
      
\medskip 
      
\noindent{\bf Catatan} Ini disebut `{\bf dekomposisi Jordan}' dari 
$\nu$.   Pernyataan dalam 231Eb adalah suatu `{\bf dekomposisi Hahn}'. 
      
      
      
\exercises{ 
\leader{231X}{Latihan dasar (a)} 
%\spheader 231Xa 
Misalkan $\Sigma$ adalah keluarga subhimpunan $A$ dari 
$\Bbb N$ sedemikian sehingga salah satu dari $A$ dan 
$\Bbb N\setminus A$ berhingga.   Tunjukkan 
bahwa $\Sigma$ adalah aljabar subhimpunan dari $\Bbb N$.   (Ini adalah 
{\bf aljabar berhingga-koberhingga} dari subhimpunan $\Bbb N$; bandingkan 
dengan 211Ra.) 
%231A 
      
\spheader 231Xb Misalkan $X$ suatu himpunan, $\Sigma$ suatu aljabar 
subhimpunan dari $X$, dan $\nu:\Sigma\to\Bbb R$ suatu fungsional aditif 
hingga.  Tunjukkan bahwa 
$\nu(E\cup F\cup G)+\nu(E\cap F)+\nu(E\cap G)+\nu(F\cap G) 
=\nu E+\nu F+\nu G+\nu(E\cap F\cap G)$ untuk semua $E$, $F$, 
$G\in\Sigma$. 
Perumumkan hasil ini untuk barisan himpunan yang lebih panjang. 
%231A 
      
\sqheader 231Xc Misalkan $\Sigma$ adalah aljabar berhingga-koberhingga dari 
subhimpunan $\Bbb N$, seperti dalam 231Xa.   Definisikan 
$\nu:\Sigma\to\Bbb Z$ dengan menetapkan 
      
\Centerline{$\nu E 
=\lim_{n\to\infty}\bigl(\#(\{i:i\le n,\,2i\in E\}) 
  -\#(\{i:i\le n,\,2i+1\in E\})\bigr)$} 
      
\noindent untuk setiap $E\in\Sigma$.   Tunjukkan bahwa $\nu$ terdefinisi 
dengan baik, aditif hingga, dan tak terbatas. 
%231Xa, 231A 
      
\spheader 231Xd Misalkan $X$ suatu himpunan dan $\Sigma$ suatu aljabar 
subhimpunan dari $X$.   (i) Tunjukkan bahwa jika 
$\nu:\Sigma\to\Bbb R$ dan $\nuprime:\Sigma\to\Bbb R$ bersifat aditif 
hingga, maka demikian pula $\nu+\nuprime$ dan $c\nu$ untuk 
sembarang $c\in\Bbb R$.   (ii) Tunjukkan bahwa jika 
$\nu:\Sigma\to\Bbb R$ bersifat aditif hingga dan $H\in\Sigma$, maka 
$\nu_H$ bersifat aditif hingga, dengan 
$\nu_H(E)=\nu(H\cap E)$ untuk setiap $E\in\Sigma$. 
%231A 
      
\spheader 231Xe Misalkan $X$ suatu himpunan, $\Sigma$ suatu aljabar 
subhimpunan dari $X$, dan $\nu:\Sigma\to\Bbb R$ suatu fungsional aditif 
hingga.   Misalkan $S$ adalah ruang linear dari fungsi-fungsi bernilai real 
pada $X$ yang dapat dinyatakan dalam bentuk $\sum_{i=0}^na_i\chi E_i$, 
dengan $E_i\in\Sigma$ untuk setiap $i$.   (i) Tunjukkan bahwa kita 
mempunyai suatu fungsional linear $\dashint:S\to\Bbb R$ yang diberikan dengan 
menuliskan 
      
\Centerline{$\dashint\sum_{i=0}^na_i\chi E_i 
=\sum_{i=0}^na_i\nu E_i$} 
      
\noindent kapan pun $a_0,\ldots,a_n\in\Bbb R$ dan 
$E_0,\ldots,E_n\in\Sigma$.   (ii) Tunjukkan bahwa jika $\nu E\ge 0$ 
untuk setiap $E\in\Sigma$, maka $\dashint f\ge 0$ kapan pun $f\in S$ dan 
$f(x)\ge 0$ untuk setiap $x\in X$.   (iii) Tunjukkan bahwa jika $\nu$ 
terbatas dan $X\ne\emptyset$, maka 
      
\Centerline{$\sup\{|\dashint f|:f\in S,\,\|f\|_{\infty}\le 1\} 
=\sup_{E,F\in\Sigma}|\nu E-\nu F|$,} 
      
\noindent dengan menuliskan 
$\|f\|_{\infty}=\sup_{x\in X}|f(x)|$. 
%231B 
      
\sqheader 231Xf Misalkan $X$ suatu himpunan, $\Sigma$ suatu aljabar-sigma 
subhimpunan dari $X$, dan $\nu:\Sigma\to \Bbb R$ suatu fungsional aditif 
hingga.   Tunjukkan bahwa pernyataan-pernyataan berikut ekuivalen: 
      
\quad (i) $\nu$ bersifat aditif terhitung; 
      
\quad (ii) $\lim_{n\to\infty}\nu E_n=0$ kapan pun $\sequencen{E_n}$ 
adalah barisan tak-meningkat di dalam $\Sigma$ dan 
$\bigcap_{n\in\Bbb N}E_n=\emptyset$; 
      
\quad (iii) $\lim_{n\to\infty}\nu E_n=0$ kapan pun 
$\sequencen{E_n}$ adalah suatu barisan di dalam $\Sigma$ dan 
$\bigcap_{n\in\Bbb N}\bigcup_{m\ge n}E_m=\emptyset$; 
      
\quad (iv) $\lim_{n\to\infty}\nu E_n=\nu E$ kapan pun 
$\sequencen{E_n}$ adalah suatu barisan di dalam $\Sigma$ dan 
      
\Centerline{$E=\bigcap_{n\in\Bbb N}\bigcup_{m\ge n}E_m 
=\bigcup_{n\in\Bbb N}\bigcap_{m\ge n}E_m$.} 
      
\noindent({\it Petunjuk}: untuk (i)$\Rightarrow$(iv), mula-mula tinjau 
$\nu$ yang nonnegatif.) 
%231D 
      
\spheader 231Xg Misalkan $X$ suatu himpunan dan $\Sigma$ suatu aljabar-sigma 
subhimpunan dari 
$X$, dan misalkan $\nu:\Sigma\to\coint{-\infty,\infty}$ suatu fungsi yang 
aditif terhitung dalam arti bahwa $\nu\emptyset=0$ dan, kapan pun 
$\sequencen{E_n}$ adalah barisan saling lepas di dalam $\Sigma$, 
$\sum_{n=0}^{\infty}\nu E_n=\lim_{n\to\infty}\sum_{i=0}^n\nu E_i$ 
terdefinisi di dalam $\coint{-\infty,\infty}$ dan sama dengan 
$\nu(\bigcup_{n\in\Bbb N}E_n)$.   Tunjukkan bahwa $\nu$ terbatas 
di atas dan mencapai batas atasnya (yakni, terdapat suatu $H\in\Sigma$ 
sedemikian sehingga $\nu H=\sup_{F\in\Sigma}\nu F$).   Selanjutnya, atau 
dengan cara lain, tunjukkan bahwa $\nu$ dapat dinyatakan sebagai selisih 
suatu ukuran berhingga total dan suatu ukuran, keduanya dengan domain 
$\Sigma$. 
%231F 
      
\leader{231Y}{Latihan lanjutan (a)} Misalkan $X$ suatu himpunan, 
$\Sigma$ suatu aljabar subhimpunan dari $X$, dan $\nu:\Sigma\to\Bbb R$ 
suatu fungsional aditif hingga yang terbatas.   Tetapkan 
      
\Centerline{$\nu^+E=\sup\{\nu F:F\in\Sigma,\,F\subseteq E\}$,} 
      
\Centerline{$\nu^-E=-\inf\{\nu F:F\in\Sigma,\,F\subseteq E\}$,} 
      
\Centerline{$|\nu|E=\sup\{\nu F_1-\nu 
F_2:F_1,\,F_2\in\Sigma,\,F_1,\,F_2\subseteq E\}$.} 
      
\noindent Tunjukkan bahwa $\nu^+$, $\nu^-$, dan $|\nu|$ semuanya 
merupakan fungsional aditif hingga yang terbatas pada $\Sigma$, serta bahwa 
$\nu=\nu^+-\nu^-$, $|\nu|=\nu^++\nu^-$.   Tunjukkan bahwa jika $\nu$ 
bersifat aditif terhitung, demikian pula $\nu^+$, $\nu^-$, dan $|\nu|$. 
($|\nu|$ kadang-kadang disebut {\bf variasi} dari $\nu$.) 
      
\header{231Yb}{\bf (b)} Misalkan $X$ suatu himpunan dan $\Sigma$ suatu 
aljabar subhimpunan dari $X$.   Misalkan $\nu_1$, $\nu_2$ adalah dua 
fungsional aditif hingga yang terbatas dan didefinisikan pada $\Sigma$. 
Tetapkan 
      
\Centerline{$(\nu_1\vee\nu_2)(E) 
=\sup\{\nu_1 F+\nu_2(E\setminus F):F\in\Sigma,\,F\subseteq E\}$,} 
      
\Centerline{$(\nu_1\wedge\nu_2)(E) 
=\inf\{\nu_1 F+\nu_2(E\setminus F):F\in\Sigma,\,F\subseteq E\}$.} 
      
\noindent Tunjukkan bahwa $\nu_1\vee\nu_2$ dan $\nu_1\wedge\nu_2$ 
merupakan fungsional aditif hingga, dan bahwa 
$\nu_1+\nu_2=\nu_1\vee\nu_2+\nu_1\wedge\nu_2$. 
Tunjukkan bahwa, dalam bahasa 231Ya, 
      
\Centerline{$\nu^+=\nu\vee 0$, 
\quad$\nu^-=(-\nu)\vee 0=-(\nu\wedge 0)$, 
\quad$|\nu|=\nu\vee(-\nu)=\nu^+\vee\nu^-=\nu^++\nu^-$,} 
      
\Centerline{$\nu_1\vee\nu_2=\nu_1+(\nu_2-\nu_1)^+$, 
\quad$\nu_1\wedge\nu_2=\nu_1-(\nu_1-\nu_2)^+$,} 
      
\noindent sehingga $\nu_1\vee\nu_2$ dan $\nu_1\wedge\nu_2$ bersifat 
aditif terhitung jika $\nu_1$ dan $\nu_2$ juga demikian. 
      
\header{231Yc}{\bf (c)} Misalkan $X$ suatu himpunan dan $\Sigma$ suatu 
aljabar subhimpunan dari $X$.   Misalkan $M$ adalah himpunan semua 
fungsional aditif hingga yang terbatas dari $\Sigma$ ke $\Bbb R$. 
Tunjukkan bahwa $M$ adalah ruang linear menurut definisi alami penjumlahan 
dan perkalian skalar.   Tunjukkan bahwa $M$ memiliki suatu urutan parsial 
$\le$ yang didefinisikan dengan mengatakan bahwa 
      
\Centerline{$\nu\le\nuprime$ jika dan hanya jika $\nu E\le\nuprime E$ untuk setiap 
$E\in\Sigma$,} 
      
\noindent dan bahwa untuk urutan parsial ini $\nu_1\vee\nu_2$, 
$\nu_1\wedge\nu_2$, sebagaimana didefinisikan dalam 231Yb, masing-masing 
adalah $\sup\{\nu_1,\nu_2\}$ dan $\inf\{\nu_1,\nu_2\}$. 
      
\header{231Yd}{\bf (d)} Misalkan $X$ suatu himpunan dan $\Sigma$ suatu 
aljabar subhimpunan dari $X$.   Misalkan $\nu_0,\ldots,\nu_n$ adalah 
fungsional-fungsional aditif hingga yang terbatas pada $\Sigma$, dan tetapkan 
      
\Centerline{$\check\nu 
E=\sup\{\sum_{i=0}^n\nu_iF_i:F_0,\ldots,F_n\in\Sigma,\,\bigcup_{i\le 
n}F_i=E,\,F_i\cap F_j=\emptyset$ untuk $i\ne j\}$,} 
      
\Centerline{$\hat\nu 
E=\inf\{\sum_{i=0}^n\nu_iF_i:F_0,\ldots,F_n\in\Sigma,\,\bigcup_{i\le 
n}F_i=E,\,F_i\cap F_j=\emptyset$ untuk $i\ne j\}$} 
      
\noindent untuk $E\in\Sigma$.   Tunjukkan bahwa $\check\nu$ dan 
$\hat\nu$ bersifat aditif hingga dan masing-masing merupakan 
$\sup\{\nu_0,\ldots,\nu_n\}$ dan 
$\inf\{\nu_0,\ldots,\nu_n\}$ di dalam himpunan terurut parsial dari 
fungsional-fungsional aditif hingga pada $\Sigma$. 
      
\header{231Ye}{\bf (e)} Misalkan $X$ suatu himpunan, $\Sigma$ suatu 
aljabar-sigma subhimpunan dari $X$; misalkan $M$ adalah himpunan terurut 
parsial dari semua fungsional aditif hingga yang terbatas dari $\Sigma$ 
ke $\Bbb R$.   (i) Tunjukkan bahwa jika $A\subseteq M$ tak kosong dan 
terbatas di atas dalam $M$, maka $A$ memiliki supremum $\check\nu$ di 
dalam $M$, yang diberikan oleh rumus 
      
$$\eqalign{\check\nu E 
=\sup\{\sum_{i=0}^n\nu_iF_i:\nu_0,\ldots,\nu_n\in 
A,\,F_0,&\ldots,F_n\in\Sigma,\,\bigcup_{i\le n}F_i=E,\cr 
&\qquad\quad F_i\cap F_j=\emptyset\text{ untuk }i\ne j\}.\cr}$$ 
      
\noindent (ii) Tunjukkan bahwa jika $A\subseteq M$ tak kosong dan 
terbatas di bawah dalam $M$, maka himpunan tersebut memiliki infimum 
$\hat\nu\in M$, yang diberikan oleh rumus 
      
$$\eqalign{\hat\nu E 
=\inf\{\sum_{i=0}^n\nu_iF_i:\nu_0,\ldots,\nu_n\in 
A,\,F_0,&\ldots,F_n\in\Sigma,\,\bigcup_{i\le n}F_i=E,\cr 
&\qquad\quad F_i\cap F_j=\emptyset\text{ untuk }i\ne j\}.\cr}$$ 
      
\header{231Yf}{\bf (f)} Misalkan $X$ suatu himpunan, $\Sigma$ suatu 
aljabar subhimpunan dari $X$, dan $\nu:\Sigma\to\Bbb R$ suatu fungsional 
aditif hingga yang nonnegatif.   Untuk $E\in\Sigma$, tetapkan 
      
\Centerline{$\nu_{ca}(E)=\inf\{\sup_{n\in\Bbb N}\nu 
F_n:\sequencen{F_n}$ adalah barisan tak-menurun di dalam $\Sigma$ dengan 
gabungan $E\}$.} 
      
\noindent Tunjukkan bahwa $\nu_{ca}$ adalah fungsional aditif terhitung 
pada $\Sigma$, dan bahwa jika $\nuprime$ adalah sembarang fungsional 
aditif terhitung dengan $\nuprime\le\nu$, maka 
$\nuprime\le\nu_{ca}$.   Tunjukkan bahwa 
$\nu_{ca}\wedge(\nu-\nu_{ca})=0$. 
      
\header{231Yg}{\bf (g)} Misalkan $X$ suatu himpunan, $\Sigma$ suatu 
aljabar subhimpunan dari $X$, dan $\nu:\Sigma\to\Bbb R$ suatu fungsional 
aditif hingga yang terbatas.   Tunjukkan bahwa $\nu$ dapat dinyatakan 
secara tunggal sebagai $\nu_{ca}+\nu_{pfa}$, dengan $\nu_{ca}$ aditif 
terhitung, $\nu_{pfa}$ aditif hingga, dan jika 
$0\le\nuprime\le|\nu_{pfa}|$ serta $\nuprime$ aditif terhitung, maka 
$\nuprime=0$. 
      
\header{231Yh}{\bf (h)} Misalkan $X$ suatu himpunan dan $\Sigma$ suatu 
aljabar subhimpunan dari $X$.   Misalkan $M$ adalah ruang linear dari 
fungsional-fungsional aditif hingga yang terbatas pada $\Sigma$, dan untuk 
$\nu\in M$, tetapkan $\|\nu\|=|\nu|(X)$, dengan mendefinisikan $|\nu|$ 
seperti dalam 231Ya.   ($\|\nu\|$ adalah {\bf variasi total} dari $\nu$.) 
Tunjukkan bahwa $\|\,\|$ adalah suatu norma pada $M$ yang menjadikan 
$M$ suatu ruang Banach.   Tunjukkan bahwa ruang fungsional-fungsional 
aditif terhitung yang terbatas pada $\Sigma$ adalah subruang linear tertutup 
dari $M$. 
      
\header{231Yi}{\bf (i)} Ulangi sebanyak mungkin hasil dalam bagian ini 
untuk fungsional bernilai kompleks. 
}%end of exercises 
      
\endnotes{ 
\Notesheader{231} Tujuan sebenarnya dari bagian ini adalah menguraikan 
dekomposisi Hahn dari suatu fungsional aditif terhitung (231E).   Uraian 
yang tidak tergesa-gesa dalam 231A-231D dimaksudkan sebagai tinjauan atas 
sifat-sifat paling elementer dari ukuran, dalam konteks `ukuran bertanda' 
yang sedikit lebih umum, dengan sifat-sifat yang hanya bersesuaian dengan 
`keaditifan' dipisahkan dari sifat-sifat yang bergantung pada `keaditifan 
terhitung'. Dalam 231Xf saya menguraikan syarat perlu dan cukup bagi suatu 
fungsional aditif hingga pada aljabar-sigma untuk bersifat aditif terhitung. 
Syarat-syarat itu dimaksudkan untuk mengisyaratkan bahwa suatu fungsional 
aditif hingga bersifat aditif terhitung jika dan hanya jika, dalam suatu 
arti, fungsional tersebut `kontinu secara sekuensial terhadap urutan'. 
Fakta bahwa 
suatu fungsional aditif terhitung dapat dinyatakan sebagai selisih 
fungsional-fungsional aditif terhitung yang nonnegatif (231F) memiliki 
padanan penting dalam teori fungsional aditif hingga: suatu 
fungsional aditif hingga dapat dinyatakan sebagai selisih 
fungsional-fungsional aditif hingga yang nonnegatif jika (dan hanya jika) 
fungsional tersebut terbatas (231Ya).   Namun, menurut saya hal ini, ataupun 
sifat-sifat lebih lanjut dari fungsional aditif hingga yang terbatas 
sebagaimana diuraikan dalam 231Xe dan 231Y, belum akan penting bagi kita 
sebelum Jilid 3. 
}%end of notes 
      
      
\discrpage 
      
